\documentclass[11pt]{article}
\usepackage[margin=0.9in]{geometry}
\usepackage{amsmath}
\usepackage{enumitem}
\usepackage{xcolor}
\setlist[itemize]{leftmargin=*,itemsep=0.2em}
\newcommand{\slide}[2]{\vspace{0.6em}\noindent\textbf{Slide #1.\ #2}\par\nopagebreak\vspace{0.1em}}

\title{Final Presentation --- Speaker Notes}
\author{Jeongwon Bae}
\date{IE 522, Spring 2026}

\begin{document}
\maketitle

\slide{1}{Title}
\begin{itemize}
\item Introduce: final presentation of our paper reproduction project.
\item We reproduced Shi et al. (2024), then engineered a pipeline that let us extend beyond the paper. I'll present the reproduction first, then the engineering, then three findings that go beyond what the paper tested.
\end{itemize}

\slide{2}{Recap: Problem and model}
\begin{itemize}
\item As shown in the midterm, this paper models emergency evacuation using three components: the community network built from OpenStreetMap, human agents with six states, and hazard agents that expand and move.
\item Since you've seen the model in the midterm, I'll move quickly.
\end{itemize}

\slide{3}{Recap: Algorithms 1 and 2}
\begin{itemize}
\item Each step, every agent runs both Algorithm 2 (human--hazard) and Algorithm 1 (human--network). Both are from the paper, unchanged in our implementation.
\end{itemize}

\slide{4}{Reproduction scope}
\begin{itemize}
\item We reproduced all three phases: 5-community network validation, Fig.\ 5 RI vs $P_\text{max} \times H_\text{max}$, and Fig.\ 6 panic sweep.
\item Then three extensions --- I'll cover those after the reproduction results.
\end{itemize}

\slide{5}{Why speed matters}
\begin{itemize}
\item A naive implementation takes about 3 minutes per single run at Pmax=2K. With 390 runs across extensions, that's on the order of hours-to-days depending on parallelism.
\item Our optimized pipeline: ~22 seconds per run, a few hours total for the whole sweep.
\item Critically: we did NOT change the algorithm. Just the functions used for path computation.
\end{itemize}

\slide{6}{Engineering 1: Profiling}
\begin{itemize}
\item The bar chart shows where time goes in a single run. The hot loop --- Algorithm 1 path recomputation --- dominates the naive case.
\item Root cause: every non-panicked agent calls networkx A* in pure Python every step.
\end{itemize}

\slide{7}{Engineering 2: scipy batch SSSP}
\begin{itemize}
\item Key insight: many agents share destinations (buildings, shelters). Compute one shortest-path tree per unique destination, reuse.
\item This converts per-agent work into per-destination work. Scipy's Dijkstra is in C, much faster than networkx in Python.
\item Measured: the Algorithm 1 section alone goes 230x faster. End-to-end including the batch SSSP overhead: 7.5x single-process.
\end{itemize}

\slide{8}{Engineering 3: A* heuristic + per-step recomputation}
\begin{itemize}
\item For single-source fallbacks (shelter redirect), we use A* with Euclidean heuristic --- faster than plain Dijkstra on spatial graphs.
\item Per-step recomputation is kept, per the paper's spec.
\item Panicked agents get noisy edge weights to model ``limited observations'' from the paper.
\item Bottom line box: algorithms unchanged, functions replaced.
\end{itemize}

\slide{9}{Engineering 4: Multiprocessing + RNG separation}
\begin{itemize}
\item 8 workers across configs --- 8x extra throughput. Network built once, shared via fork copy-on-write.
\item RNG separation is subtler but critical: 3 independent streams for humans, hazards, sim dynamics. Without it, changing Pmax changes the hazards too --- apples-to-oranges.
\end{itemize}

\slide{10}{Engineering 5: Benchmark}
\begin{itemize}
\item We ran a head-to-head: naive Pmax=2K run vs optimized.
\item 7.5x single-process, 60x with 8 workers.
\item What this enabled: 10 seeds per config, 5 communities, three extensions --- without taking shortcuts.
\end{itemize}

\slide{11}{E: PSU-UP animation}
\begin{itemize}
\item This replaces the toy animation from the midterm.
\item Configuration: paper's baseline (Pmax=2K, Hmax=5, epsilon\_p=10\%).
\item Watch: hazard appears, impacted agents reroute to shelter, some are casualties, the rest arrive.
\item Outcome corresponds to RI $\approx$ 36\% which we'll see on slide 14.
\end{itemize}

\slide{12}{Phase 1: Network validation}
\begin{itemize}
\item Building counts match the paper's Table V within 0--2\% across all 5 communities.
\item Edge counts: PSU-UP, UVA-C, VT-B are close. RA-PA and KOP-PA are way off --- paper reports E/N ratios of 6.5--6.8, unreproducible with any OSMnx config. Treated as a paper data anomaly.
\end{itemize}

\slide{13}{Phase 1: PSU-UP figures}
\begin{itemize}
\item Fig.\ 3 reproduction: walk network with buildings in yellow, shelters in green.
\item Fig.\ 4 reproduction: pedestrian flow snapshots at five timesteps during a hazard scenario.
\end{itemize}

\slide{14}{Phase 2a: RI vs Pmax x Hmax}
\begin{itemize}
\item Fig.\ 5 reproduction. Our RI is 2--4\%p higher than paper uniformly.
\item Critical confirmation: RI is independent of Pmax across 2K--8K ($\Delta < 0.2$\%p). This is the paper's core finding, and we reproduced it.
\end{itemize}

\slide{15}{Phase 2b: RS/RC/RL vs epsilon\_p}
\begin{itemize}
\item Fig.\ 6 reproduction with 9 epsilon\_p points instead of 5.
\item All three monotonic trends reproduced. Near-exact match at boundary points.
\item Paper's main conclusion holds: panic is the dominant driver of survival/casualty outcomes.
\end{itemize}

\slide{16}{Differences and unresolvable gaps}
\begin{itemize}
\item Four categories of remaining gap: all traced to unpublished paper data.
\item Shelter list not in the paper --- we swept shelter count and found 62\% of buildings best matches RS at epsilon\_p=90\%.
\item Casualty formula not specified --- we use a distance-weighted probability.
\item Hazard seed not specified --- we swept 200 seeds and picked 1135.
\end{itemize}

\slide{17}{Extension overview}
\begin{itemize}
\item Three questions the paper didn't answer, which our fast pipeline let us tackle.
\item Move quickly --- details on the next three slides.
\end{itemize}

\slide{18}{Extension 1: Pmax breakpoint}
\begin{itemize}
\item Extended Pmax from 2K to 97K across three Hmax levels.
\item Independence holds up to 20K. Breaks at 50K by +8.5\%p.
\item Mechanism: 969 buildings at 20 occupancy is about 20K capacity. At 50K we overflow --- agents queue outdoors and get hit by hazards.
\end{itemize}

\slide{19}{Extension 2: Hmax saturation}
\begin{itemize}
\item Extended Hmax from 5 to 30 at Pmax=2K.
\item RI saturates at $\sim$86\% --- more hazards don't push past this ceiling.
\item Likely $\sim$14\% of the network is structurally unreachable by hazards.
\end{itemize}

\slide{20}{Extension 3: Finer epsilon\_p}
\begin{itemize}
\item 9 points vs paper's 5.
\item Curves are smooth, monotonic, diminishing marginal --- no phase transitions.
\item Implication for policy: panic management has highest payoff at low baseline panic.
\end{itemize}

\slide{21}{Comparison alpha: Pmax animation}
\begin{itemize}
\item Same hazard scenario, different Pmax. Watch for outdoor queuing on the right.
\end{itemize}

\slide{22}{Comparison beta: epsilon\_p animation}
\begin{itemize}
\item Same hazard, different panic. Notice how the right panel's agents wander randomly instead of moving toward shelters.
\end{itemize}

\slide{23}{Comparison gamma: Hmax animation}
\begin{itemize}
\item Hazard counts differ by design --- say this explicitly. Left has 5 hazard events, right has 30.
\item Visual takeaway: right panel has most of the network covered in red, agents hit repeatedly.
\end{itemize}

\slide{24}{Summary}
\begin{itemize}
\item Reproduction successful within 2--4\%p across all paper targets.
\item Three novel findings: Pmax breakpoint, Hmax saturation, smooth panic effect.
\item Engineering enabled all of this without algorithm changes.
\end{itemize}

\slide{25}{Future work}
\begin{itemize}
\item Paper's own: response agents, social-force, MDP panic model.
\item Our own gap analysis suggests a panic-dependent casualty formula could close the RC ceiling gap.
\end{itemize}

\slide{26}{References}
\begin{itemize}
\item Acknowledge OSMnx, scipy. Code available on GitHub.
\item Thank audience, invite questions.
\end{itemize}

\end{document}
